\documentclass[10pt,a4paper]{amsart}
\usepackage[margin=19mm]{geometry}
\usepackage{amsmath,amssymb,mathtools}
\usepackage[scr=rsfs,cal=euler]{mathalfa}
\usepackage{xfrac}
\usepackage[unicode,hidelinks,bookmarks=false]{hyperref}
\newcommand{\ie}{i.e.\ }
\newcommand{\cf}{cf.\ }
\newcommand{\resp}{respectively\ }
\newcommand{\Subsection}{\subsection}
\providecommand{\nobreakdash}{\nobreak\hbox{-}}
\setlength{\emergencystretch}{2em}
\multlinegap0pt
\usepackage{amssymb}
\usepackage{xcolor}
\usepackage[shortlabels]{enumitem}
\usepackage{mathtools}
\newtheorem{conjecture}{Conjecture}
\newtheorem{theorem}{Theorem}
\newcommand{\K}{\mathcal{K}}
\newcommand{\R}{\mathbb{R}}
\newcommand{\Z}{\mathbb{Z}}
\newcommand{\floor}[1]{\lfloor #1 \rfloor}
\title{Quantitative Stability of the Betke-Henk-Wills Conjecture}
\author{Chao Wang}
\date{Source: arXiv:2602.06662v2 (CC BY 4.0)}
\begin{document}
\maketitle

\noindent\emph{Source and licence.} Quantitative Stability of the Betke-Henk-Wills Conjecture. Version arXiv:2602.06662v2 (CC BY 4.0), distributed by arXiv under the Creative Commons Attribution 4.0 International licence, http://creativecommons.org/licenses/by/4.0/, as recorded in the arXiv licence metadata for this exact version and shown on the abstract page https://arxiv.org/abs/2602.06662v2. Full licence notice: \texttt{licenses/arxiv-2602.06662-CC-BY-4.0.txt}.\par\smallskip

\noindent\textit{Editorial scope.} Counted unit: the article's theorem thm:box (source0.tex lines 152--154). The article's complete original proof of it is delivered with it (source0.tex lines 156--174, one \texttt{proof} environment of 19 lines). No mathematics is written, repaired or re-derived: the statement and the proof are the article's own lines, in the article's own notation, and no retained line is substituted or re-flowed.  Retained uncounted as the source-local context the proof needs: lines 92--97.  Nothing was omitted from any retained span, so the package carries no omission record.  No retained line contains a citation, so the wrapper carries no bibliography and no bibliography entry.  No retained line contains a figure environment, an image-inclusion command or drawing code, so no image is delivered and the package declares no asset dependency.  Editorial formatting only, in addition: an AMS \texttt{amsart} preamble that restates the article's own declarations and macros used by the retained lines, namely the environments \texttt{conjecture}, \texttt{theorem} on separate counters, together with the standard packages those lines need.  The retained environments print in this document as Conjecture 1, Theorem 1 in order of appearance, whereas the article prints the same items with the numbering of its own full document.  The complete native source of the article as distributed in the arXiv e-print for this exact version is bundled unchanged as \texttt{sources/194112/source0.tex}.
\begin{conjecture}[Betke-Henk-Wills] \label{conj:main}
For any $K \in \K^d_0$ and $\Lambda \in \mathcal{L}^d$:
\begin{equation}
G(K, \Lambda) \le \prod_{i=1}^d \left\lfloor \frac{2}{\lambda_i(K, \Lambda)} + 1 \right\rfloor.
\end{equation}
\end{conjecture}

\begin{theorem}[Conjecture for Boxes] \label{thm:box}
Let $K = \{x \in \R^d : |x_i| \le \alpha_i \}$ be an axis-aligned box with semi-axes $\alpha_1 \ge \dots \ge \alpha_d > 0$. Then Conjecture \ref{conj:main} holds.
\end{theorem}

\begin{proof}
For the standard lattice $\Z^d$, the successive minima are achieved by the standard basis vectors $e_i$, yielding:
\begin{equation}
\lambda_i(K, \Z^d) = \frac{1}{\alpha_i}.
\end{equation}
The lattice point enumerator factorizes due to the Cartesian product structure:
\begin{equation}
G(K, \Z^d) = \prod_{i=1}^d G([-\alpha_i, \alpha_i], \Z) = \prod_{i=1}^d (2\floor{\alpha_i} + 1).
\end{equation}
Let $\mathcal{R}(K)$ denote the right-hand side of the inequality. Substituting $\lambda_i = 1/\alpha_i$:
\begin{equation}
\mathcal{R}(K) = \prod_{i=1}^d \floor{\frac{2}{\lambda_i} + 1} = \prod_{i=1}^d \floor{2\alpha_i + 1}.
\end{equation}
It suffices to show that for any real $x \ge 0$, $2\floor{x} + 1 \le \floor{2x + 1}$.
Let $x = I + f$ where $I \in \Z_{\ge 0}$ and $0 \le f < 1$.
The LHS is $2I + 1$. The RHS is $\floor{2I + 2f + 1} = 2I + \floor{2f + 1}$.
Since $f \ge 0$, $\floor{2f + 1} \ge 1$. Thus, $2I + 1 \le 2I + \floor{2f+1}$ holds.
Multiplying over all dimensions confirms the conjecture.
\end{proof}

\end{document}
